\subsection{Isometric linkage of the original cutoff models}
\label{sec:geometry-original-metric-link}

We next prove the original two-cutoff linkage on the
\texttt{RootOrigin}/\texttt{GoodOrigin} data. The argument
uses one common cover and one identity of regularizing
potentials, rather than comparing two arbitrary choices
of metric constants.

\begin{theorem}[Common model and exact two-level power-sum comparison]
\label{thm:geometry-original-metric-link}
Retain the actual fixed data and the original two cutoff
models with correction sets \(S_0\subset S\).
Use the source's hypotheses that the normal projective
models have the full prescribed generic torsion, their
corrected cotangent modules and labelled images are
locally free at correction primes, and their actual
patched flags are flags of vector bundles. Assume the
\texttt{RootOrigin} identifications with the same fixed
data, including the cubical pullbacks, the actual root
line presentations, the original singular Hodge norms
and the original regularizations. At the common primes
the full levels agree. All ramified coefficient primes
and \(2\) are in \(S_0\).
Then the two models have a common normal projective
generically finite comparison model on which the actual
root pairs have vertical difference supported at \(S\setminus S_0\)
and have exactly isometric generic smooth metrics for
all sufficiently small positive regularization parameters.
The comparison retains the actual normalization degrees.
On this common model the relative Chern identities
(RC4)--(RC5) and the original fibre power-sum vanishing
hold at the new primes. Consequently the original
\texttt{Prop6} comparison of normalized power-sum intersections
holds whenever the prescribed limiting intersections exist.
No extra height term arises from making the comparison model
or choosing its projective coordinate presentation.
\end{theorem}
\begin{proof}
\paragraph{One common geometric model.}
Let \(K_f\) denote the fixed generic flag field in
\texttt{FixedData}. The embeddings \(g_i:K_f\to K_i\)
of the two origin fields identify the same map to the
original moduli component and the same generic flag:
these are exactly \texttt{g\_over} and \texttt{g\_flag}.
Choose the two finite extensions in a common algebraic
closure of \(K_f\), and take their compositum \(K_*\).
No linear disjointness is required. The integers
\[
                 e_i=[K_*:K_i],\qquad d_i=[K_i:K_f]
\]
satisfy \(e_0d_0=e_1d_1=[K_*:K_f]\).
Take the closure of the common generic correspondence
in the product of the two projective models, then its
normalization in \(K_*\). These are proper constructions;
normalization is finite since the models are excellent.
They give a normal projective integral model \(X_*\)
with dominant maps to both original models of degrees
\(e_i\). In particular it is flat over \(\mathbf Z\).
The two induced morphisms to the original moduli model
coincide: they agree generically and the target is separated.
The original \texttt{SABMap} is a group pullback and
pulls back the coefficient action and cubical line.
Thus the two semiabelian families, their polarizations
and their generic tangent identifications on \(X_*\)
are literally identified by base change from this common
moduli map. This argument does not use an isogeny between
unrelated abelian varieties.

The ordinary cotangent bundle of a smooth group commutes
with arbitrary base change. At a correction prime the
module of character differentials is the inverse image
of the span of the specified full-level character
differentials in \(\omega/p^{R_p}\omega\).
All generic torsion points already occur over the original
field, so the full list after pullback is exactly the same
list, not a proper subset. Their differentials commute with
dominant base change by the original cubical pairing
(or by Theorem~\ref{thm:geometry-canonical-chart-order}
and its cotangent identification at level \(p^{R_p}\)).
In the original projected cotangent images this description
gives their generated pullbacks. Since those images are
locally free, taking their duals commutes with pullback.
For the \texttt{htW} label the original inverse-different
factor is an invertible ideal from the coefficient ring;
it pulls back as that same factor. This proves exact
compatibility also at the ramified primes of \(S_0\).
For \texttt{dR} at the primes outside \(S_0\), the
coefficient algebra is unramified and split after the
prescribed embedding of the constants. Its eigenspaces are
direct summands and commute with pullback.
These are the actual formulas \texttt{htOmEig},
\texttt{htDual}, \texttt{htW} and \texttt{Lat.patch}, with
the prescribed two different conventions for the two labels.

For a locally free lattice with a flag by subbundles,
pullback of its flag is the actual intersection with the
pulled-back rational flag: locally the subbundle has a
locally free complement, and that complement is torsion
free after pullback to the integral comparison model.
This proves the identification without commuting an
arbitrary unsaturated intersection with base change.
The generic flags of the two origins agree by
\texttt{g\_flag}. Consequently their pulled-back root
lines are the actual graded lines of the two patched
lattices, with the prescribed reversal and duality for
the second block. At every point outside the new
correction fibres those lattices, flags and graded lines
agree, including their natural generic identifications.
At a new prime the old and new flags are respectively
the de Rham and the corrected flag of the same group.
The canonical root-pair isomorphism off the new fibres
is therefore a rational section of the old/new line ratio
with divisor supported on the new fibres, not an
arbitrary rational section with a horizontal divisor.

For applying the relative Chern identities at a new prime,
work over the local ring of the appropriate constant
number field at that prime, which is a DVR. Its ring of
integers embeds in the normal comparison charts by
integral closedness. We may include all the constants
and the prescribed roots of unity from the finite levels
in \(K_*\), already present in the full-level data.
In particular \(w\) has \(w^{p-1}=up\) at the new prime.
At each such prime all coefficient labels are unramified;
the determinant and flag bundles needed in the relative
comparison are locally free by the original hypotheses.
If a further auxiliary refinement is needed for the
images in the relative Chern diagram, use
Lemma~\ref{lem:relative-chern-lattice-diagram}.
By Corollary~\ref{cor:geometry-original-chart-identity},
the actual chart test identities hold at its height-one
fields. Thus Theorem~\ref{thm:relative-chern-two-identities}
gives (RC4)--(RC5) for the two pulled-back flags.
Its left side uses the tangent roots in the first block
and the reversed dual roots in the second block, exactly
the \texttt{rootIdx}/\texttt{rootSgn} convention; the
different is a unit at this prime. The old Lie power
sums on the right are therefore the source's
\texttt{liePs}, with the indicated signs in the
second block. For the new flag the exact sequence (HF1)
gives the fibre power-sum vanishing, including
multiplicities when intersecting on a nonreduced fibre.
All these statements keep the original roots by pullback.

The permutation used here is independent of the auxiliary
constant-field place: since the coefficient field \(E\)
is abelian Galois, at an unramified rational prime the
arithmetic Frobenius element \(\mathrm{Fr}_p\) is independent
of the prime of \(E\). Precomposition of the labelled
embeddings by that element has the source convention
\(\chi_{\phi\tau}=\chi_\tau^p\). It either preserves
both CM half-types or interchanges them, according to its
restriction to the original imaginary quadratic field.
This is the original Frobenius shape, not a new choice
of orientation.

\paragraph{Exact metric agreement.}
On the common complex interior the old and new singular
Hodge metrics on a given generic root line coincide,
under the actual generic identification. Indeed the
Hermitian form of the polarized abelian scheme is
functorial for the cubical pullbacks just described.
A flag quotient norm is the corresponding quotient of
the same positive Gram form (the square root of the
successive leading-minor ratio); taking the dual changes
it to the inverse norm in precisely the source's second
block. At complex points the correction primes are
invertible, so the patched lattices do not change this
generic flag or the metric. This is the original
\texttt{Flag.IsHodgeLine}/\texttt{SecNorm} prescription.

For clarity we verify the actual regularization, not just
the singular metric. Write \(z_i\) for the projections of
a common complex point in the original coordinate models.
For each model the map to the common analytic space
\(\mathrm{Sys}.X\) is identified by the fixed homogeneous
coordinate block \(\mathrm{Fx}.yf\).
The two blocks are proportional at every point of the
comparison: their projective morphisms coincide on the
generic fibre and hence everywhere by separatedness.
The factor relating two nonzero homogeneous
representatives is nonzero. The original \texttt{pt\_inj}
and \texttt{\(\pi\)\_pt} show that the corresponding
analytic points are the same, say \(x\).
One may make this verification on the dense analytic
open covered by the original realizations; their
\texttt{J} is the top ideal and their interior clause
selects the common abelian locus. There is no assertion
that a chosen analytic lift of an arbitrary point is
unique, only that its image in the fixed analytic space
is the same.

Let \(h^{\rm sing}_{i,\tau k}\) and \(h^t_{i,\tau k}\)
be the positive homogeneous metric functions for the
singular and smoothed lines, with the original inverse
norm convention of \texttt{HermLine.h}.
For each prefix \(j\) the three defining identities
\texttt{fund\_sing}, \texttt{Re.line\_eq}, and
\texttt{fund\_reg} imply on the complex interior
\[
 \sum_{k\le j}\log h^t_{i,\tau k}(z_i)
  =\sum_{k\le j}\log h^{\rm sing}_{i,\tau k}(z_i)
      +\mathrm{Sys.ph}(\tau,j,x)
      -\mathrm{Sys.phT}(\tau,j,t,x),\quad 0<t\le t_0.       \tag{GL1}
\]
This follows by substituting the equality for the
smoothed reference metric into the prefix identity.
Taking consecutive prefix differences (with empty prefix
zero) shows that for each root, the correction of its
logarithmic metric is a fixed function of \(t,x\), identical
for the two origins. The singular isometry therefore
implies the exact smoothed isometry. In particular the
argument uses the full scalar norm in
\texttt{SecNorm}, not just the curvature of the metric.

To express it in the original homogeneous presentation,
add a very ample coordinate block to embed \(X_*\) as
a closed subscheme of the product of the two coordinate
spaces and this additional projective space. Pullback of
the old coordinate blocks defines the required block maps.
A rational section of the old/new root ratio is represented
by a quotient of two multihomogeneous forms \(f_r/g_r\),
with multidegree difference the old degree minus the new
degree. To see existence, trivialize the ratio over a
dense coordinate chart by the corresponding coordinate
monomial, express the remaining rational function as a
quotient in the additional projective coordinate chart,
and homogenize, multiplying by powers of the chosen
coordinate entries to achieve the specified degree
difference. Clear scalar denominators equally in numerator
and denominator. If \(e_0,e_1\) are the two original
root trivializing scalars, choose the generic ratio
\(e_1/e_0\), as in \texttt{Flag.TwoLevels}.
This is precisely the rational section of the old/new
ratio specified above. The defining equality
\texttt{SecNorm} for a common generic section gives
\[
              |f_r(z)|\,h^t_{1,r}(z)
                        =|g_r(z)|\,h^t_{0,r}(z).
                                                               \tag{GL2}
\]
For example on a coordinate chart the norm of an
element \(s\) of the root line is the absolute value of
\(s/(e_i\prod z_j^{a_{i,j}})\) times the homogeneous
coordinate scale and divided by \(h^t_{i,r}\).
Cancel the common section and use the generic equality
\(f_r/g_r=e_1/e_0\) with the stated degrees.
This fixes the orientation of (GL2), rather than
deducing it from equality of curvatures.
The equality extends to all complex points by continuity
of the smooth positive metrics and the homogeneous forms.
The common generic line isomorphism has no horizontal
zero or pole; equivalently it is a smooth isometry on
the entire complex fibre. If the source requires a
positive-parameter metric beyond the range \(t\le t_0\),
extend the second metric by transport of the first metric
using this isomorphism. For \(0<t\le t_0\) keep the
original two metrics. The extension does not affect any
limit at zero and is smooth and conjugation invariant
at each fixed parameter. This gives \texttt{IsChange}
with the literal vertical and isometry clauses.

The auxiliary ample line also agrees on the common model:
its coordinate block defines the same morphism to the
fixed projective space, and its metric is the same
\(\mathrm{Fx.hf}\). The factor \(k_A\) is the same by
\texttt{kA\_eq}. Thus this isometry does not introduce a
change in the perturbing ample line or a metric
renormalization.

\paragraph{Normalized intersections and the local sum.}
In the original normalization \(\delta_i=[c.F:\mathbf Q]d_i\).
On the common model take \(\delta_*=[c.F:\mathbf Q][K_*:K_f]\).
Thus
\[
                       \delta_*=e_i\delta_i .
\]
For every smooth arithmetic intersection of the original
pulled-back lines the projection formula multiplies the
unnormalized number by \(e_i\), and the geometric
intersection by the same factor. Division by \(\delta_*\)
therefore gives exactly the original normalized number,
with \(k_A\) unchanged. In particular for each fixed
positive smoothing parameter in the original range, the
intersection functions of the common model are the two
original ones. Taking their prescribed limits preserves
this equality; no convergence uniform in the cutoffs or
the generic degree is needed. Any auxiliary birational
modification used for the local lattices has degree one
in this comparison. This proves the \texttt{PP} and
geometric parts of the original linkage.

Apply the vertical comparison lemma of
Section~\ref{sec:geometry-two-level-weighted-bridge} on
this common model before taking the smoothing limit.
The root \(p_{k,\tau}\) is the sum of the \(k\)-th
powers of the indicated full flag roots, and \(Q\) is
any homogeneous complementary polynomial of degree \(d+1-k\).
All mixed terms involving \(Q(x)-Q(y)\) vanish by the
new fibre power-sum identity, as in that lemma.
The remaining terms are exactly (RC4)--(RC5) with
\(Q\) on the old roots. Since the metrics are isometric
generically, these are purely vertical arithmetic
intersections with their original \(\log p\) factors.
The local class is \(k\) times the indicated old
\((k-1)\)-st power sum times the principal fibre
\([\operatorname{div}(p)]\), with the Frobenius orbit
coefficients. Equivalently it is the original
\(\mathrm{bracket}_{j,\phi_p,p,k,\tau}\,Q\) evaluated on
the old fibre, including the fibre's multiplicities.

Specialization of a proper flat model identifies the
degree of the intersection of line classes on this fibre
with their generic geometric intersection degree.
This follows, for example, by writing the classes as
differences of relatively very ample classes and
polarizing the constancy of the generic and special
Hilbert polynomial; hence it applies to these signed
root monomials as well. Over a constant extension the
ramification multiplicity in \(\operatorname{div}(p)\)
and the residue degree in the arithmetic degree are both
included; working over \(\mathbf Z\) adds all these
places automatically. Thus after division by
\(\delta_*\) the sum for each \(p\) is
\(\log p\) times the old normalized geometric functional
on the bracket polynomial. There is no loss proportional
to the chosen normalization degree.

In the notation of the source, writing \(B_0,B\) for the
cutoffs, \(a\) for the original component and
\(\operatorname{geo}_{B_0,a}\) for its normalized
geometric functional, we have proved
\[
 \begin{split}
 &\mathrm{PPc}(B_0,a,p_{k,\tau}Q)
       -\mathrm{PPc}(B,a,p_{k,\tau}Q)\\
 &\hspace{10mm}=\sum_{p\in S\setminus S_0}
       \log p\,\operatorname{geo}_{B_0,a}
                (\mathrm{bracket}_{j,\phi_p,p,k,\tau}Q).
 \end{split}                                                     \tag{GL3}
\]
If the original proportionality
\(\operatorname{geo}_{B_0,a}=v_a\deg_{\mathrm{Fl}}\)
is used, this is exactly the source's \texttt{Prop6}:
the finite orbit formula (RC5) equals its geometric-series
coefficient, so the summand is \(v_a\) times its
\texttt{deltaTerm}. The set of primes is exactly
\(\{p:B_0<p\le B,\ p\ {\rm prime},\ p\notin\mathrm{bad}\}\)
when the original cutoff sets are used.
This proves the required two-level comparison for the
actual root models, with the same normalization and
regularization; it makes no claim of small complexity for
the common comparison cover.
\end{proof}

The theorem is an ordinary arithmetic-intersection
reconstruction of \texttt{link\_step} and
\texttt{LinkedTower} under the original model, metric and
limit inputs, using the new local lattice suppliers.
In particular it establishes the two-level power-sum
formula needed for the weighted height calculation.
It does not by itself prove the source's remaining
same-tower flag-height, remainder, proportionality,
volume and finite section-rank estimates; nor does it
replace proof of existence of the original
\texttt{GoodOrigin} data or convergence by a formal
comparison theorem. The conditional weighted-height
bound in Section~\ref{sec:geometry-two-level-weighted-bridge}
now has the exact local correction and linkage supplier.
